\section{Anomalies de liens (étudiant 1)}
\label{sec:liens}

\subsection{Protocole et modèles}

Faute d'étiquette « lien faux », nous ajoutons au graphe $r\,|E|$ liens entre des paires non reliées tirées uniformément, $r\in\{1,5,10,20\}\,\%$, \emph{avant} le découpage des liens (\texttt{RandomLinkSplit}, 10\,\% validation, 20\,\% test). Les modèles sont entraînés par reconstruction, sans étiquette : les liens injectés sont dans le graphe d'entraînement. L'encodeur ne voit que les liens d'entraînement (70\,\%, injectés compris), et chaque lien du graphe contaminé reçoit le score $s_{ij}=1-\sigma(z_i^\top z_j)$ (moyenne $\mu$ pour le VGAE). Pour les baselines, le score est l'opposé de l'indice (Jaccard, Adamic--Adar) ou la distance en km.

Les entrées sont lat, lon et population, standardisées. Le GAE a deux couches GCN ($3\to64\to32$) séparées par une ReLU, et la perte est l'entropie croisée sur les liens d'entraînement et des paires négatives tirées au hasard. Le VGAE a une première couche GCN identique, puis deux couches GCN en parallèle pour $\mu$ et $\log\sigma$ ; sa perte vaut $\mathcal{L}_{\mathrm{rec}}+\frac{1}{N}\mathrm{KL}$ ($N$ nœuds). Adam (taux $0{,}01$, 200 époques) sous PyTorch Geometric~\cite{fey2019pyg}, graines $42, 123, 456, 789, 2026$. Nous reportons l'AUC et la précision moyenne (AP) ; la précision@10 est trop instable (écart-type jusqu'à $0{,}25$).

\subsection{Résultats}

\begin{table}[t]
\centering
\footnotesize
\caption{Liens injectés : moyenne sur 5 graines, l'écart-type est entre parenthèses (dernier chiffre : $0{,}952(4)=0{,}952\pm0{,}004$). Meilleure valeur en gras. Une AP aléatoire vaut $r/(1+r)$, soit $0{,}010$, $0{,}048$, $0{,}091$ et $0{,}167$.}
\label{tab:links}
\setlength{\tabcolsep}{3.5pt}
\begin{tabular}{@{}lcccccccc@{}}
\toprule
& \multicolumn{2}{c}{$r=1\,\%$} & \multicolumn{2}{c}{$r=5\,\%$} & \multicolumn{2}{c}{$r=10\,\%$} & \multicolumn{2}{c}{$r=20\,\%$} \\
\cmidrule(lr){2-3}\cmidrule(lr){4-5}\cmidrule(lr){6-7}\cmidrule(lr){8-9}
Méthode & AUC & AP & AUC & AP & AUC & AP & AUC & AP \\
\midrule
Distance (km) & $0{,}914(12)$ & $0{,}080(7)$ & $0{,}916(5)$ & $0{,}298(9)$ & $0{,}915(3)$ & $0{,}452(6)$ & $0{,}916(3)$ & $0{,}613(7)$ \\
Jaccard & $0{,}949(6)$ & $0{,}098(3)$ & $0{,}947(2)$ & $0{,}349(3)$ & $0{,}946(2)$ & $0{,}516(4)$ & $0{,}946(1)$ & $0{,}679(2)$ \\
Adamic--Adar & $\mathbf{0{,}952}(4)$ & $0{,}098(3)$ & $\mathbf{0{,}951}(2)$ & $0{,}351(3)$ & $\mathbf{0{,}951}(1)$ & $0{,}519(3)$ & $\mathbf{0{,}950}(1)$ & $\mathbf{0{,}683}(2)$ \\
\midrule
GAE & $0{,}945(7)$ & $0{,}190(25)$ & $0{,}936(4)$ & $0{,}429(19)$ & $0{,}929(4)$ & $\mathbf{0{,}552}(8)$ & $0{,}917(3)$ & $0{,}650(7)$ \\
VGAE & $0{,}943(9)$ & $\mathbf{0{,}193}(14)$ & $0{,}933(3)$ & $\mathbf{0{,}434}(15)$ & $0{,}927(4)$ & $0{,}551(10)$ & $0{,}916(2)$ & $0{,}654(5)$ \\
\bottomrule
\end{tabular}
\end{table}

Adamic--Adar a le meilleur AUC à tous les taux ($0{,}95$ contre $0{,}917$ à $0{,}945$), sans apprentissage ni attribut (tableau~\ref{tab:links}). Les autoencodeurs gardent l'avantage en AP quand les anomalies sont rares ($0{,}19$ contre $0{,}10$ à $r=1\,\%$), puis il s'inverse à $r=20\,\%$ ($0{,}68$ contre $0{,}65$). La distance seule donne déjà $0{,}915$ : des paires tirées au hasard sont presque toujours éloignées.

GAE et VGAE sont à moins de $0{,}005$ l'un de l'autre, comme l'écart-type ; au test de Student apparié, $p>0{,}2$ dans cinq cas sur huit et $p\approx0{,}03$--$0{,}05$ dans les trois autres. Rien ne montre que le VGAE améliore le GAE, nous gardons le plus simple.

Sur le graphe d'origine (graine 42), quatre des dix liens les mieux classés sont communs aux deux modèles ; la plupart relient un grand aéroport à une destination éloignée (Paris--Cayenne, Londres--La Barbade) (atypiques, mais probablement réels).
% TODO : vérifier ces liaisons sur une source externe (OpenFlights, Wikipédia) et ajuster la phrase.
Un cas est une vraie erreur : le VGAE classe cinquième Ottawa--Iqaluit, or le nœud « Ottawa » a pour coordonnées $(-27{,}9;\,144{,}6)$ et pour pays l'Australie, alors que ses 26 voisins sont presque tous canadiens ou américains. Le lien est correct, ce sont les attributs du nœud qui sont faux.

\subsection{Ablation du GAE}

\begin{table}[t]
\centering
\footnotesize
\caption{Ablation du GAE ($r=10\,\%$, 5 graines, écart-type entre parenthèses). Référence : 2 couches, attributs standardisés, $\dim Z=32$.}
\label{tab:ablation}
\begin{tabular}{@{}lcc@{}}
\toprule
Variante & AUC & AP \\
\midrule
GAE complet & $0{,}929(4)$ & $0{,}552(8)$ \\
\midrule
1 couche GCN & $0{,}846(10)$ & $0{,}365(14)$ \\
Sans standardisation & $0{,}500(0)$ & $0{,}091(0)$ \\
Sans attributs ($X=\mathbf 1$) & $0{,}812(3)$ & $0{,}202(4)$ \\
$\dim Z=8$ (au lieu de 32) & $0{,}930(3)$ & $0{,}558(16)$ \\
\bottomrule
\end{tabular}
\end{table}

Avec une seule couche GCN, l'AUC passe de $0{,}93$ à $0{,}85$ (tableau~\ref{tab:ablation}). Sans standardisation il vaut exactement $0{,}5$ et l'AP est celle du hasard ($0{,}091$) : la population, de l'ordre de $10^7$, fait probablement saturer le sigmoïde et tous les scores sont égaux. Sans attributs, l'AUC reste à $0{,}81$ mais l'AP tombe à $0{,}20$. Une dimension latente de 8 donne le même résultat que 32, qui n'est donc pas nécessaire.

Limites : les liens injectés sont tirés au hasard et donc faciles à reconnaître ; un AUC de $0{,}93$ ici ne dit pas que les vraies erreurs seront trouvées. Le top-10 vient d'une seule graine.
